\section{Details of Non-Implications (Non-Additive Valuations)}
\label{sec:cex-nonadd-extra}

\begin{table*}[!htb]
\centering
\caption{Non-implications among fairness notions (non-additive valuations).}
\label{table:cex-nonadd}
\footnotesize
\begin{tabular}{cccccccc}
\toprule & \tiny valuation & \tiny marginals & \tiny identical & \tiny $n$ & \tiny entitlements & &
\\\midrule \hfill\makecell[l]{EF+GAPS $\nfimplies$ \\ PROP1 or PROPm}\hfill
    & supermod & $\ge 0$ bival\textsuperscript{\ref{foot:nonneg-bival}}
    & yes & $n \ge 2$ & equal & \cref{cex:ef-not-prop-supmod} & trivial
\\[\defaultaddspace] PPROP $\nfimplies$ PROP1
    & supermod & $\ge 0$ bival\textsuperscript{\ref{foot:nonneg-bival}}
    & yes & $n \ge 3$ & equal & \cref{cex:ef-not-prop-supmod} & trivial
%
\\\midrule PROP $\nfimplies$ M1S\textsuperscript{\ref{foot:unit-dem}}
    & unit-dem & $\ge 0$ & yes & $n=2$ & equal
    & \cref{cex:prop-not-m1s-unit-demand} & \textbf{new}
\\[\defaultaddspace] PROP1 $\nfimplies$ PROPm\textsuperscript{\ref{foot:unit-dem}}
    & unit-dem & $\ge 0$ & yes & $n=2$ & equal
    & \cref{cex:ud:prop1-not-propm} & \textbf{new}
\\[\defaultaddspace] PROP $\nfimplies$ PPROP\textsuperscript{\ref{foot:unit-dem}}
    & unit-dem & $\ge 0$ & yes & $n=3$ & equal
    & \cref{cex:ud:prop-not-pprop} & \textbf{new}
\\[\defaultaddspace] PROPm $\nfimplies$ PROPavg\textsuperscript{\ref{foot:unit-dem}}
    & unit-dem & $\ge 0$ & yes & $n=3$ & equal
    & \cref{cex:ud:propm-not-propavg} & \textbf{new}
\\[\defaultaddspace] GAPS $\nfimplies$ PROPx
    & unit-dem & $\ge 0$ & yes & $n=3$ & equal & \cref{cex:gaps-not-propx} & \textbf{new}
\\[\defaultaddspace] MMS $\nfimplies$ EF1 or PROPm
    & unit-dem & $\ge 0$ & yes & $n=4$ & equal
    & \cref{cex:ud:mms-not-ef1-propm} & \textbf{new}
%
\\\midrule EF $\nfimplies$ MMS
    & submod & $\ge 0$ bival\textsuperscript{\ref{foot:nonneg-bival}}
    & yes & $n=2$ & equal & \cref{cex:ef-not-mms-pmrf} & folklore
\\[\defaultaddspace] EEF $\nfimplies$ EF1 or PPROP
    & submod & $\ge 0$ bival\textsuperscript{\ref{foot:nonneg-bival}}
    & yes & $n=3$ & equal & \cref{cex:eef-not-ef1-pprop-pmrf} & \textbf{new}
\\[\defaultaddspace] MEFS $\nfimplies$ EF1
    & submod & $\ge 0$ bival\textsuperscript{\ref{foot:nonneg-bival}}
    & yes & $n=2$ & equal & \cref{cex:mefs-not-ef1-pmrf} & \textbf{new}
\\[\defaultaddspace] PROP $\nfimplies$ M1S
    & submod & $\ge 0$ bival\textsuperscript{\ref{foot:nonneg-bival}}
    & yes & $n=2$ & equal & \cref{cex:prop-not-m1s-uniform-matroid} & \textbf{new}
\\[\defaultaddspace] MMS $\nfimplies$ APS
    & submod & $> 0$ & yes & $n=2$ & equal & \cref{cex:mms-not-aps-n2-submod} & \cite{babaioff2023fair}
\\[\defaultaddspace] GAPS+PPROP $\nfimplies$ PROP1
    & subadd & $\ge 0$ bival\textsuperscript{\ref{foot:nonneg-bival}}
    & yes & $n=3$ & equal & \cref{cex:gaps-not-prop1-subadd} & \textbf{new}
%
\\\midrule EF1 $\nfimplies$ MXS
    & submod & $\{0, 1\}$ & yes & $n=2$ & equal & \cref{cex:ef1-not-mxs-pmrf} & \textbf{new}
\\[\defaultaddspace] GAPS $\nfimplies$ PROP1 or EF1
    & subadd & $\{0, 1\}$ & yes & $n=2$ & equal & \cref{cex:gaps-not-ef1-prop1-subadd} & \textbf{new}
\\[\defaultaddspace] GAPS+EF1 $\nfimplies$ PROP1
    & subadd & $\{0, 1\}$ & yes & $n=2$ & equal & \cref{cex:ef1-not-prop1-subadd} & \textbf{new}
\\[\defaultaddspace] \hfill\makecell[l]{PAPS+PPROP \\ $\nfimplies$ PROPm or M1S}\hfill
    & subadd & $\{0, 1\}$ & yes & $n=3$ & equal & \cref{cex:paps-pprop-not-propm-binary-subadd} & \textbf{new}
\\[\defaultaddspace] GMMS $\nfimplies$ APS
    & subadd & $\{0, 1\}$ & yes & $n=3$ & equal & \cref{cex:gmms-not-aps-binary-subadd} & \textbf{new}
\\[\defaultaddspace] PROPavg $\nfimplies$ PROPx
    & supermod & $\{0, 1\}$ & yes & $n \ge 3$ & equal & \cref{cex:ef-not-prop-supmod} & \textbf{new}
\\ \bottomrule
\end{tabular}

\begin{tightenum}
\item[*] \label{foot:unit-dem}Result also holds for
    submodular cancelable valuations with positive marginals.
\item[\textdagger] \label{foot:nonneg-bival}Result holds for both
    $\{0, 1\}$ marginals and positive bivalued marginals.
\end{tightenum}
\end{table*}


\subsection{Supermodular Valuations}
\label{sec:cex-extra:supmod}

\begin{lemma}[EF+GAPS $\nfimplies$ PROP1]
\label[lemma]{cex:ef-not-prop-supmod}
Let $0 \le \eps < 1/27$ and $k \in \mathbb{N}$. For any $x \ge 0$, let
\[ f(x) \defeq x\eps + \max(0, x-k). \]
Consider a fair division instance with $n$ equally-entitled agents
having an identical valuation function $v$ over $4n$ goods,
where $v(X) = f(|X|)$ for all $X \subseteq [4n]$. Then $v$ is supermodular.

Let $A$ be the allocation where each agent gets 4 goods. Then $A$ is EF+GAPS+GMMS.
If $k = 5$, then no PROP1 or PROPm allocation exists.
If $n \ge 3$ and $k = 8$, then $A$ is PPROP but no PROP1 or PROPx allocation exists.
Additionally if $\eps = 0$, then $A$ is PROPavg.
\end{lemma}
\begin{proof}
For any $X \subseteq [4n]$ and $g \in [4n] \setminus X$, we have
\[ v(g \mid X) = \begin{cases}
\eps & \textrm{ if } |X| \le k
\\ 1 + \eps & \textrm{ if } |X| \ge k
\end{cases}. \]
Since $v(g \mid X)$ is non-decreasing in $|X|$, $v$ is supermodular.

Since all agents have the same bundle, $A$ is EF and GMMS.
Set the price of each good to 1 to get that $A$ is GAPS.

The proportional share is $4-k/n+4\eps$.
For $k = 5$, the proportional share is at least $3/2$,
but $f(5) = 5\eps$ and $f(6) = 1+6\eps < 3/2$.
In any allocation, some agent will have at most 4 goods,
and that agent would not be PROP1 or PROPm satisfied.

If $n \ge 3$ and $k = 8$, then the proportional share is at least $4/3$.
Since $f(8) = 8\eps$ and $f(4) = 4\eps$, $A$ is pairwise PROP.
In any allocation, some agent gets at most 4 goods,
and $f(5) = 5\eps < 4/3$, so that agent is not PROP1-satisfied,
and $f(9) = 1 + 9\eps < 4/3$, so that agent is not PROPx-satisfied.
Additionally, if $\eps = 0$, for any two agents $i$ and $j$,
we have $v(A_j \mid A_i) = f(8) - f(4) = 0$. Hence, $A$ is PROPavg.
\end{proof}

\subsection{Unit-Demand Valuations}
\label{sec:cex-extra:unit-demand}

It is known that unit-demand functions are submodular and cancelable.
We first show how to perturb a unit-demand function such that
submodularity and cancelability remain preserved and all marginals become positive.

\begin{lemma}
\label[lemma]{thm:ud-perturb}
Let $u: 2^{[m]} \to \mathbb{R}_{\ge 0}$ be a function and $\eps \in \mathbb{R}_{\ge 0}$.
For any $X \subseteq [m]$, define $\uhat(X) \defeq u(X) + |X|\eps$.
If $u$ is submodular, then $\uhat$ is also submodular.
If $m \le 3$ and $u$ is cancelable, then $\uhat$ is also cancelable.
\end{lemma}
\begin{proof}
If $u$ is submodular, then for any $S, T \subseteq [m]$, we get
\[ \uhat(S) + \uhat(T) - \uhat(S \cup T) - \uhat(S \cap T)
= u(S) + u(T) - u(S \cup T) - u(S \cap T) \ge 0. \]
Hence, $\uhat$ is also submodular.

Let $u$ be cancelable.
For every $T \subseteq [m]$ and $S_1, S_2 \subseteq [m] \setminus T$,
if $\uhat(S_1 \cup T) > \uhat(S_2 \cup T)$,
we need to show that $\uhat(S_1) > \uhat(S_2)$
to prove that $\uhat$ is cancelable.
This is trivial when $T = \emptyset$ or $\eps = 0$.

If $|S_1| = |S_2|$, then $u(S_1 \cup T) > u(S_2 \cup T)$.
Since $u$ is cancelable, we get $u(S_1) > u(S_2)$, which implies $\uhat(S_1) > \uhat(S_2)$.

Let $\eps > 0$. If $S_1 = \emptyset$, we get
$\uhat(S_1 \cup T) > \uhat(S_2 \cup T) \implies u(S_2 \mid T) < - |S_2|\eps$,
which is a contradition since $u$ is non-negative. Hence, $S_1 \neq \emptyset$.
If $S_2 = \emptyset$, we get $\uhat(S_1) \ge |S_1|\eps > 0 = \uhat(S_2)$.
When $m \le 3$, we have considered all cases.
Hence, if $m \le 3$ and $u$ is cancelable, then $\uhat$ is also cancelable.
\end{proof}

\begin{lemma}[PROP $\nfimplies$ M1S]
\label[lemma]{cex:prop-not-m1s-unit-demand}
Consider a unit-demand function $v$ over 3 goods such that $v(1) = v(2) = 4$ and $v(3) = 3$.
Let $0 \le \eps < 1/6$ and $\vhat(X) \defeq v(X) + 2|X|\eps$ for all $X \subseteq [3]$.
Consider a fair division instance with 3 goods and
2 equally-entitled agents having valuation function $\vhat$.
Then the allocation $A = (\{1, 2\}, \{3\})$ is PROP but not M1S.
\end{lemma}
\begin{proof}
$A$ is PROP since $\vhat(A_1) = 4+4\eps$, $\vhat(A_2) = 3+2\eps$, and the PROP share is $2+3\eps$.
Suppose $A$ is M1S and agent 2's M1S certificate for $A$ is $B$.
$B_2 \neq \emptyset$, since $B$ is EF1-fair to agent 2.
Moreover, $\vhat(B_2) \le \vhat(A_2) = 3 + 2\eps$, so $B_2 = \{3\}$. Hence, $B_1 = \{1, 2\}$.
However, $\min_{g \in B_1} \vhat(B_1 \setminus \{g\}) = 4 + 4\eps > 3 + 2\eps = \vhat(B_2)$.
Hence, agent 2 is not EF1-satisfied by $B$, which contradicts the fact that $B$
is agent 2's M1S certificate for $A$. Hence, $A$ is not M1S.
\end{proof}

\begin{example}[PROP1 $\nfimplies$ PROPm]
\label[example]{cex:ud:prop1-not-propm}
Consider a unit-demand function $v$ over 2 goods such that $v(1) = 30$ and $v(2) = 10$.
Let $0 \le \eps < 1$ and $\vhat(X) \defeq v(X) + |X|\eps$ for all $X \subseteq [2]$.
Consider a fair division instance with 2 goods and
2 equally-entitled agents having valuation function $\vhat$.
Then the allocation $A \defeq (\{1, 2\}, \emptyset)$ is PROP1 but not PROPm.
\end{example}

\begin{example}[PROP $\nfimplies$ PPROP]
\label[example]{cex:ud:prop-not-pprop}
Consider a unit-demand function $v$ over 3 goods such that $v(1) = 30$ and $v(2) = v(3) = 11$.
Let $0 \le \eps < 1$ and $\vhat(X) \defeq v(X) + |X|\eps$ for all $X \subseteq [3]$.
Consider a fair division instance with 3 goods and
3 equally-entitled agents having valuation function $\vhat$.
Then the allocation $A \defeq (\{1\}, \{2\}, \{3\})$ is PROP but no alloction is PPROP.
\end{example}

\begin{example}[MMS $\nfimplies$ EF1, PROPm]
\label[example]{cex:ud:mms-not-ef1-propm}
Consider a unit-demand function $v$ over 5 goods such that
$v(1) = 200$ and $v(2) = 30$ and $v(3) = v(4) = v(5) = 10$.
Let $0 \le \eps < 1$ and $\vhat(X) \defeq v(X) + |X|\eps$ for all $X \subseteq [5]$.
Consider a fair division instance with 5 goods and
4 equally-entitled agents having valuation function $\vhat$.
Then the allocation $A \defeq (\{1, 2\}, \{3\}, \{4\}, \{5\})$
is MMS but not EF1 and not PROPm.
(The maximin share is $10 + \eps$. The proportional share is $50 + (5/4)\eps$.)
\end{example}

\begin{example}[PROPm $\nfimplies$ PROPavg]
\label[example]{cex:ud:propm-not-propavg}
Consider a unit-demand function $v$ over 3 goods such that
$v(1) = 75$, $v(2) = 30$, and $v(3) = 10$.
Let $0 \le \eps < 1$ and $\vhat(X) \defeq v(X) + |X|\eps$ for all $X \subseteq [3]$.
Consider a fair division instance with 3 goods and
3 equally-entitled agents having valuation function $\vhat$.
Then the allocation $A \defeq (\{1, 3\}, \{2\}, \emptyset)$ is PROPm but not PROPavg.
(The proportional share is $25+\eps$.)
\end{example}

\subsection{Submodular Valuations}
\label{sec:cex-extra:submod}

\begin{definition}[PMRF]
\label[definition]{defn:pmrf}
Let $P = (P_1, \ldots, P_k)$ be a partition of a finite set $M$.
Define $\PMRF(P)$ as the function $f: 2^M \to [|M|]$ where
$f(X) \defeq \sum_{i=1}^k \boolone(X \cap P_i \neq \emptyset)$.
(Here $\boolone(C)$ is 1 if condition $C$ is true and 0 otherwise.)
For any $\eps \in \mathbb{R}_{\ge 0}$, define $\PMRF_{+\eps}(P)$
as the function $f$ where $f(X) = \PMRF(P)(X) + |X|\eps$.
\end{definition}

\begin{lemma}
\label[lemma]{thm:pmrf-submod}
For any partition $P$ over a finite set $M$,
$f \defeq \PMRF_{+\eps}(P)$ is submodular, and $f(g \mid X) \in \{\eps, 1+\eps\}$
for any $X \subseteq M$ and $g \in M \setminus X$.
\end{lemma}
\begin{proof}[Proof sketch]
$\PMRF(P)$ is the rank function of a partition matroid, and hence, is submodular.
% TODO: prove.
The submodularity of $f$ follows from \cref{thm:ud-perturb}.
\end{proof}

\begin{example}[EF $\nfimplies$ MMS]
\label[example]{cex:ef-not-mms-pmrf}
Let $P \defeq (\{r_1, r_2\}, \{g_1, g_2\})$, $\eps \in \mathbb{R}_{\ge 0}$,
and $v \defeq \PMRF_{+\eps}(P)$.
Consider a fair division instance with 2 equally-entitled agents
having identical valuation function $v$.
Then the maximin share is at least $2 + 2\eps$,
as exemplified by the partition $(\{r_1, g_1\}, \{r_2, g_2\})$.
The allocation $P$ is EF but not MMS,
since each agent's value for her own bundle is $1 + 2\eps$.
\end{example}

\begin{example}[EF1 $\nfimplies$ MXS]
\label[example]{cex:ef1-not-mxs-pmrf}
Let $P \defeq (\{r_1, r_2\}, \{g_1, g_2\}, \{b\})$ and $v \defeq \PMRF(P)$.
Consider a fair division instance with 2 equally-entitled agents
having identical valuation function $v$.
Then the allocation $A \defeq (\{r_1, r_2\}, \{g_1, g_2, b\})$ is EF1
but agent 1 is not MXS-satisfied by $A$.
\end{example}

\begin{example}[MEFS $\nfimplies$ EF1]
\label[example]{cex:mefs-not-ef1-pmrf}
Let $P \defeq (\{a\}, \{b\}, \{c\}, \{d\}, \{e_1, e_2\}, \{f_1, f_2\}, \{g_1, g_2\})$.
Let $v \defeq \PMRF_{+\eps}(P)$ for $0 \le \eps \le 1/2$.
Consider a fair division instance with 2 equally-entitled agents
having identical valuation function $v$.
Let $A \defeq (\{a, e_1, f_1, g_1\}, \{b, c, d, e_2, f_2, g_2\})$.
Since $v(A_2) = 6 + 6\eps$ and $v(A_1) = 4 + 4\eps$,
agent 2 is envy-free in $A$ but agent 1 is not EF1-satisfied by $A$.
However, agent 1 is MEFS-satisfied by $A$, and her MEFS-certificate for $A$ is
$B \defeq (\{a, b, c, d\}, \{e_1, e_2, f_1, f_2, g_1, g_2\})$,
since $v(B_1) = v(A_1) = 4 + 4\eps$ and $v(B_2) = 3 + 6\eps$.
\end{example}

\begin{example}
\label[example]{cex:eef-not-ef1-pprop-pmrf}
Let $M \defeq \{a_j\}_{j=1}^6 \cup \{b_j\}_{j=1}^6 \cup \{c_j\}_{j=1}^4$.
Let $P$ be a partition of $M$ where $X \in P$ iff $X = \{a_j, b_j\}$ for some $j \in [6]$
or $X = \{c_j\}$ for some $j \in [4]$.
Let $v \defeq \PMRF_{+\eps}(P)$ for $0 \le \eps < 1/6$.
Consider a fair division instance with 3 equally-entitled agents
having identical valuation function $v$.

Let $A \defeq (\{a_j\}_{j=1}^6, \{b_j\}_{j=1}^6, \{c_j\}_{j=1}^4)$.
Then $v(A_1) = v(A_2) = 6(1+\eps)$ and $v(A_3) = 4(1+\eps)$.
Then $A$ is not EF1 and not PPROP.
But $A$ is EEF since agents 1 and 2 are envy-free,
and agent 3's EEF-certificate is $B$ where
$B_1 \defeq \{a_1, a_2, a_3, b_1, b_2, b_3\}$,
$B_2 \defeq \{a_4, a_5, a_6, b_4, b_5, b_6\}$, and $B_3 = A_3$.
Then $v(B_3) = 4(1+\eps)$ and $v(B_1) = v(B_2) = 3 + 6\eps < 4$.
\end{example}

\begin{example}[PROP $\nfimplies$ M1S]
\label[example]{cex:prop-not-m1s-uniform-matroid}
Let $v: 2^{[10]} \to \mathbb{R}_{\ge 0}$ be a function defined as $v(X) \defeq \min(6, |X|)$.
Then $v$ is submodular (since $v$ is the rank function of a uniform matroid)
and $v$ has binary marginals (i.e., $v(g \mid X) \in \{0, 1\}$
for all $X \subseteq [10]$ and $g \in [m] \setminus X$).
Let $0 \le \eps < 1$ and $\vhat(X) \defeq v(X) + |X|\eps$.

Consider a fair division instance with 2 equally-entitled agents
having identical valuation function $\vhat$.
Then $A \defeq ([4], [10] \setminus [4])$ is PROP but agent 1 is not M1S-satisfied.
\end{example}

\begin{lemma}[MMS $\nfimplies$ APS for $n=2$, Remark 2 of \cite{babaioff2023fair}]
\label[lemma]{cex:mms-not-aps-n2-submod}
Let $C \defeq \{\{1, 2, 3\}, \{1, 5, 6\}, \{2, 4, 6\}, \{3, 4, 5\}\}$.
Let $v: 2^{[6]} \to \mathbb{R}_{\ge 0}$ be a function defined as
\[ v(X) \defeq \begin{cases}
0 & \text{ if } X = \emptyset
\\ 2 & \text{ if } |X| = 1
\\ 4 & \text{ if } |X| = 2
\\ 5 & \text{ if } |X| = 3 \text{ and } X \not\in C
\\ 6 & \text{ if } X \in C \text{ or } |X| \ge 4
\end{cases}. \]
Let $0 \le \eps < 1/3$ and define $\vhat$ as $\vhat(X) \defeq v(X) + |X|\eps$.
Then $\vhat$ is submodular.

Consider a fair division instance with 2 equally-entitled agents
having identical valuation function $\vhat$.
The maximin share is $5 + 3\eps$, but the APS is $6 + 3\eps$.
Moreover, there exists an MMS allocation but no allocation is APS.
\end{lemma}
\begin{proof}
For any $X \subseteq [6]$, and $g \in [6] \setminus X$, we have
\begin{tightenum}
\item $|X| \le 1 \implies \vhat(g \mid X) = 2 + \eps$.
\item $|X| = 2 \implies \vhat(g \mid X) \in \{1 + \eps, 2 + \eps\}$.
\item $|X| = 3 \implies \vhat(g \mid X) \in \{\eps, 1 + \eps\}$.
\item $|X| = 4 \implies \vhat(g \mid X) = \eps$.
\end{tightenum}
Hence, $\vhat(g \mid X)$ decreases with $|X|$, so $\vhat$ is submodular.

The MMS is less than 6 because for every set in $C$,
its complement is not in $C$.
The APS is $6 + 3\eps$ by the dual definition of APS (\cref{defn:aps-dual}),
by setting $x_S = 1/4$ when $S \in C$.

For identical valuations, an MMS allocation always exists.
For identical valuations, in any allocation, some agent gets a bundle of value at most the MMS.
Since the APS is strictly greater than the MMS, that agent is not APS-satisfied.
Hence, no allocation is APS.
\end{proof}

\subsection{Subadditive Valuations}
\label{sec:cex-extra:subadd}

\begin{lemma}
\label[lemma]{thm:floor-ceil-value}
For any $k \in \mathbb{N}$, define the set functions $f_1, f_2: 2^{[m]} \to \mathbb{R}$ as
$f_1(S) = \floor{|S|/k}$ and $f_2(S) = \ceil{|S|/k}$, respectively.
Then $f_1$ is superadditive, $f_2$ is subadditive,
and for any $S \subseteq [m]$, $g \in [m] \setminus S$, and $i \in [2]$,
we have $f_i(g \mid S) \in \{0, 1\}$.
\end{lemma}

\begin{example}
\label[example]{cex:gaps-not-ef1-prop1-subadd}
Let $([2], [9], (v_i)_{i=1}^2, \eqEnt)$ be a fair division instance where
$v_i(S) \defeq \ceil{|S|/4}$ for all $i \in [2]$ and $S \subseteq [9]$.
Then $v_i$ is subaditive for $i \in [2]$ by \cref{thm:floor-ceil-value}.
Let $A$ be an allocation where agent 1 has 3 goods and agent 2 has 6 goods.
Then $A$ is GAPS+EFX+PROPx but not EF1 or PROP1.

Each agent's APS is $\le 1$, because if we price each good at $1/9$,
then a bundle is affordable iff its value is at most 1.
Each agent's proportional share is $3/2$.
Note that $A$ is not \EFXZero{} (c.f.~\cref{defn:efx0-goods});
the difference between EFX and \EFXZero{} is crucial here.
\end{example}

\begin{example}
\label[example]{cex:ef1-not-prop1-subadd}
Let $([2], [7], (v_i)_{i=1}^2, \eqEnt)$ be a fair division instance where
$v_i(S) \defeq \ceil{|S|/4}$ for all $i \in [2]$ and $S \subseteq [7]$.
Then $v_i$ is subaditive for $i \in [2]$ by \cref{thm:floor-ceil-value}.
Let $A$ be an allocation where agent 1 has 2 goods and agent 2 has 5 goods.
Then $A$ is GAPS+EF1 but not PROP1.
(Note that in our definition of PROP1 (c.f.~\cref{defn:prop1}), we require
$v_i(A_i \cup \{g\}) > v_i([m])/2$ for some $g \in [m] \setminus A_i$,
not just $v_i(A_i \cup \{g\}) \ge v_i([m])/2$.
This nuance is crucial to this counterexample.)
\end{example}

\begin{lemma}[optimal bin packing is subadditive]
\label[lemma]{thm:bin-packing}
Let $(s_j)_{j=1}^m$ be a sequence where $0 \le s_j \le 1$ for all $j \in [m]$
($s_j$ is called the \emph{size} of item $j$).
For any $X \subseteq [m]$, define $\optBP_s(X)$ as the smallest integer $k$
such that a partition $(P_1, \ldots, P_k)$ of $X$ exists
such that $\sum_{j \in P_i} s_j \le 1$ for all $i \in [k]$.
Then $\optBP_s$ is subadditive and has binary marginals,
i.e., $\optBP_s(X \cup \{g\}) - \optBP_s(X) \in \{0, 1\}$.
\end{lemma}
\begin{proof}
Pick any two disjoint sets $X, Y \subseteq [m]$.
Let $k_X \defeq \optBP_s(X)$ and $k_Y \defeq \optBP_s(Y)$.
Let $(P_1, \ldots, P_{k_X})$ and $(Q_1, \ldots, Q_{k_Y})$ be the corresponding partitions.
Then $(P_1, \ldots, P_{k_X}, Q_1, \ldots, Q_{k_Y})$ is a feasible partition of $X \cup Y$,
so $\optBP_s(X \cup Y) \le k_X + k_Y$.
\end{proof}

\begin{example}
\label[example]{cex:gaps-not-prop1-subadd}
Let $s = (1/5, 1/5, 3/5, 3/5, 3/5, 3/5)$. Let $0 \le \eps < 1/3$.
Consider a fair division instance with 3 equally-entitled agents
having an identical valuation function $v$ over 6 goods.
$v(X) \defeq \optBP_s(X) + |X|\eps$
($\optBP_s$ is defined in \cref{thm:bin-packing}).
Then by \cref{thm:bin-packing}, $v$ is subadditive,
and for every $X \subseteq [6]$ and $g \in [6] \setminus X$,
we have $v(g \mid X) \in \{\eps, 1+\eps\}$.

Then allocation $A = (\{1, 2\}, \{3, 4\}, \{5, 6\})$ is EFX, pairwise PROP, and groupwise APS
(set the price of the first two goods to $1/2$ and the last four goods to 1).
However, agent 1 is not PROP1-satisfied in $A$, since $v([6]) = 4 + 6\eps$,
and $v(A_1 \cup \{g\}) = 1 + 3\eps$ for all $g \in [6] \setminus A_1$

(Note that unlike \cref{cex:ef1-not-prop1-subadd}, this example
works even for the weaker definition of PROP1.)
\end{example}

\begin{lemma}[PAPS+PPROP $\nfimplies$ PROPm or M1S for binary subadd]
\label[lemma]{cex:paps-pprop-not-propm-binary-subadd}
Let $A_1$, $A_2$, and $A_3$ be pairwise-disjoint sets where
$A_1$ contains 7 goods of size $0.6$, $A_2$ contains 21 goods of size $0.4$,
and $A_3$ contains 21 goods of size $0.4$. Let $M \defeq A_1 \cup A_2 \cup A_3$.
Let $v(X) \defeq \optBP(X)$ for all $X \subseteq M$.
Consider a fair division instance with 3 equally-entitled agents
having the same valuation function $v$ over $M$.
Then allocation $A \defeq (A_1, A_2, A_3)$ is EFX+PPROP+PAPS,
but $A$ is not PROPm-fair or M1S-fair to agent 1.
\end{lemma}
\begin{proof}
First, we show that $v(S) = \max(|S \cap A_1|, \ceil{|S|/2})$ for any $S \subseteq M$.
To see this, suppose $S$ contains $a$ items of size $0.6$ and $b$ items of size $0.4$.
If $b \le a$, we require $a$ to pack the items and $\ceil{(a+b)/2} \le a$.
If $b \ge a$, then we require $\ceil{(a+b)/2}$ bins to pack the items and $a \le (a+b)/2$.
As a corollary, we get that for all $S \subseteq M \setminus A_1$,
we have $v(S \mid A_1) = \max(0, \ceil{(|S|-|A_1|)/2})$.

$v(A_2) = v(A_3) = 11$ and $v(A_1) = 7$, so agents 2 and 3 are envy-free.
For any $S \subseteq A_2$, $v(S \mid A_1) > 0 \iff |S| \ge 8$,
and $|S| \ge 8 \implies |A_2 \setminus S| \le 13 \implies v(A_2 \setminus S) \le 7 = v(A_1)$.
Hence, agent 1 doesn't EFX-envy agent 2. Similarly, she doesn't EFX-envy agent 3.
Hence, $A$ is EFX.

$v(A_2 \cup A_3) = 21$, and $v(A_1 \cup A_2) = v(A_1 \cup A_3) = 14$.
Hence, $A$ is pairwise PROP. We will now show that $A$ is pairwise APS.
Set the price of each item to 1.
It's easy to see that the allocation restricted to agents 2 and 3 is APS.
Now suppose we restrict the allocation to agents 1 and 2.
Then $S \subseteq A_1 \cup A_2$ is affordable iff $|S| \le |A_1 \cup A_2|/2 = 14$.
For $|S| = 14$, we get $v(S) = 7$, so the APS is at most 7.
Hence, $A$ is pairwise APS satisfied.

$v(M)/3 = 25/3 = 8 + 1/3$.
For $S \subseteq A_2$ such that $|S| = 8$, we get that $v(S \mid A_1) = 1$,
but $v(A_1 \cup S) = 8 < v(M)/3$. Hence, $A$ is not PROPm-fair to agent 1.
Also, $v(A_1 \cup \{g\}) = v(A_1) = 7$ for all $g \in M \setminus A_1$,
so $A$ is also not PROP1-fair to $A$.

Suppose $A$ is M1S-fair to agent 1. Let $B$ be her M1S-certificate for $A$.
Then $v(B_1) \le v(A_1) = 7$ and $B$ is EF1-fair to agent 1.
$v(B_1) \le 7$ implies $|B_1| \le 14$. Since $|M| = 49$, we get that
$\max(|B_2|, |B_3|) \ge \ceil{|M - B_1|/2} \ge 18$.
Hence, some agent $i \in \{2, 3\}$ has a bundle of value at least $9$.
If one good is removed from their bundle, its value can drop by at most $8$,
so $B$ is not EF1-fair to agent 1, which is a contradiction.
Hence, $A$ is not M1S-fair to agent 1.
\end{proof}

\begin{lemma}[GMMS $\nfimplies$ APS]
\label[lemma]{cex:gmms-not-aps-binary-subadd}
Let there be 15 items of sizes
$65/96$, $31/96$, $31/96$, $31/96$, $23/96$, $23/96$, $23/96$, $17/96$,
$11/96$, $7/96$, $7/96$, $7/96$, $5/96$, $5/96$, $5/96$.
Consider a fair division instance with 3 equally-entitled agents having
the same valuation function $v \defeq \optBP_s$ over the 15 items (c.f.~\cref{thm:bin-packing}).
Then $v$ is subadditive and has binary marginals.
Also, the AnyPrice share is at least 2, but the maximin share is 1.
Hence, no APS allocation exists, but the leximin allocation is GMMS.
\end{lemma}
\begin{proof}
The total size of the items is $3 \times 97/96$.
By Lemma C.1 of \cite{babaioff2023fair}, no $3$-partition of the items exists
such that the total size of each partition is at least $97/96$.
Hence, in every 3-partition $M$, some bundle has total size at most 1, so that bundle has value at most 1.
So, the maximin share is 1.

By Lemma C.1 of \cite{babaioff2023fair}, there exist 6 sets $(S_j)_{j=1}^6$ of items
such that each set has total size $97/96$ (and hence value 2) and each item appears in exactly two sets.
Hence, the APS is at least 2 (c.f.~\cref{defn:aps-dual}).
\end{proof}
